% CP-VI-0084
% target_chapter: VI/32 — Tangent Spaces and Local Geometry
% solution_provenance: FRESH_CANONICAL_AUTHORED_SOLUTION

\subsection*{CP-VI-0084: Tangent space to a plane curve}

Let $X=V(F)\subseteq\mathbb A_k^2$ and $p=(a,b)\in X$. Show that the Zariski tangent space is
\[
T_pX=\{(u,v)\in k^2:F_x(p)u+F_y(p)v=0\}.
\]
Apply this to the cusp $y^2=x^3$ at the origin.

\medskip
\noindent\textbf{Solution.}

A tangent vector is a $k[\varepsilon]/(\varepsilon^2)$-point reducing to $p$. Substituting $(a+\varepsilon u,b+\varepsilon v)$ into $F$ gives
\[
F(p)+\varepsilon(F_x(p)u+F_y(p)v),
\]
so the coefficient of $\varepsilon$ must vanish. For $F=y^2-x^3$ at $(0,0)$ both partial derivatives vanish, hence $T_0X=k^2$. Its tangent-space dimension exceeds the curve dimension one, detecting the singularity.
